\chapter{Lesson 8: Directions and Transportation}

This lesson covers getting around: location words (in front of, behind, next to), asking
for and giving directions, and talking about transportation. The grammar points are the
instrumental/directional particle \ko{(으)로}, the range expressions \ko{에서 \til 까지}
and \ko{부터 \til 까지}, and how long something takes with \ko{걸리다}.

\section{Vocabulary}

\begin{table}[H]
\centering
\begin{tabular}{lll}
\toprule
Korean & RR & English \\
\midrule
\ko{버스} & beoseu & bus \\
\ko{지하철} & jihacheol & subway \\
\ko{택시} & taeksi & taxi \\
\ko{기차} & gicha & train \\
\ko{비행기} & bihaenggi & airplane \\
\ko{걸어서} & georeoseo & on foot \\
\ko{정류장} & jeongnyujang & bus stop \\
\ko{역} & yeok & station \\
\ko{공항} & gonghang & airport \\
\ko{타다} & tada & to get on, ride (takes \ko{을/를}) \\
\ko{내리다} & naerida & to get off (takes \ko{에서}) \\
\ko{갈아타다} & garatada & to transfer \\
\ko{걸리다} & geollida & to take (time) \\
\ko{길} & gil & road, way \\
\ko{사거리} & sageori & intersection \\
\ko{신호등} & sinhodeung & traffic light \\
\ko{오른쪽 / 왼쪽} & oreunjjok / oenjjok & right / left \\
\ko{쭉 / 똑바로} & jjuk / ttokbaro & straight \\
\ko{건너다} & geonneoda & to cross \\
\ko{돌다} & dolda & to turn \\
\ko{가깝다 / 멀다} & gakkapda / meolda & to be near / far \\
\bottomrule
\end{tabular}
\end{table}

\section{Location words}

Position relative to something is expressed as \emph{reference noun} + \emph{position
noun} + \ko{에} (for existence) or \ko{에서} (for actions): \ko{은행 앞에} (in front of the
bank), \ko{책상 위에} (on the desk).

\begin{center}
\begin{tikzpicture}[every node/.style={font=\small}]
  \draw[thick, fill=gray!15] (0,0) rectangle (2.6,1.8);
  \node at (1.3,0.9) {\ko{안} an};
  \node[font=\scriptsize] at (1.3,0.45) {inside};
  \node[above] at (1.3,1.85) {\ko{위} wi (above, on)};
  \node[below] at (1.3,-0.05) {\ko{아래 / 밑} arae / mit (below)};
  \node[left, align=center] at (-0.1,0.9) {\ko{옆} yeop\\ \scriptsize beside};
  \node[right, align=center] at (2.7,0.9) {\ko{옆} yeop\\ \scriptsize beside};
  % top view: a building seen from above, street at the bottom
  \draw[thick, fill=gray!15] (5.2,0.85) rectangle (7.8,1.65);
  \node[font=\scriptsize] at (6.5,1.25) {\ko{은행} (top view)};
  \fill[gray!25] (4.6,-0.3) rectangle (8.4,0.1);
  \node[font=\scriptsize] at (8.9,-0.15) {street};
  \node[font=\small] at (6.5,0.45) {\ko{앞} ap (front)};
  \node[font=\small] at (6.5,2.0) {\ko{뒤} dwi (behind)};
  \node[align=left, anchor=west, font=\scriptsize] at (9.6,0.95) {also \ko{밖} bak (outside),\\ \ko{사이} sai (between),\\ \ko{건너편} geonneopyeon\\ (across from)};
\end{tikzpicture}
\end{center}

\begin{example}[Where is it?]
\ko{어디에 있어요?} Eodie isseoyo? (Where is it?)
\begin{itemize}
  \item \ko{우체국은 은행 옆에 있어요.} Uchegugeun eunhaeng yeope isseoyo. (The post
        office is next to the bank.)
  \item \ko{편의점은 학교 건너편에 있어요.} Pyeonuijeomeun hakgyo geonneopyeone isseoyo.
        (The convenience store is across from the school.)
  \item \ko{도서관하고 식당 사이에 커피숍이 있어요.} Doseogwanhago sikdang saie keopisyobi
        isseoyo. (There's a coffee shop between the library and the cafeteria.)
  \item \ko{정류장 앞에서 기다릴게요.} Jeongnyujang apeseo gidarilgeyo. (I'll wait in front
        of the bus stop.) --- action: \ko{에서}.
\end{itemize}
\end{example}

\section{Grammar 1: \ko{(으)로} (by means of; toward)}

\ko{(으)로} attaches to nouns. After a vowel or the consonant \ko{ㄹ} it is \ko{로};
after other consonants it is \ko{으로}.

\begin{enumerate}
  \item \textbf{Means / instrument}: \ko{버스로} beoseuro (by bus), \ko{지하철로}
        jihacheollo (by subway, \ko{ㄹ} + \ko{로}), \ko{한국어로} hangugeoro (in Korean),
        \ko{젓가락으로} jeotgarageuro (with chopsticks).
  \item \textbf{Direction} (toward): \ko{오른쪽으로} oreunjjogeuro (to the right),
        \ko{학교로} hakgyoro (toward school), \ko{이쪽으로 오세요} (come this way).
\end{enumerate}

\begin{remark}
\ko{에} and \ko{(으)로} can both mark a destination. \ko{에} names the end point
(\ko{학교에 가요}, I go to school), while \ko{(으)로} focuses on direction (\ko{학교로
가요}, I head toward school). With turning verbs and \ko{쪽} (side), \ko{(으)로} is
required: \ko{왼쪽으로 도세요} (turn left).
\end{remark}

\begin{example}[Means of transport]
Getting to campus:
\begin{itemize}
  \item \ko{학교에 어떻게 와요? --- 버스로 와요.} Hakgyoe eotteoke wayo? --- Beoseuro
        wayo. (How do you get to school? --- By bus.)
  \item \ko{저는 걸어서 다녀요.} Jeoneun georeoseo danyeoyo. (I walk; lit.\ go back and
        forth by walking.) --- ``on foot'' is \ko{걸어서}, not \ko{*발로}.
  \item \ko{시청역에서 1호선으로 갈아타세요.} Sicheongyeogeseo ilhoseoneuro garataseyo.
        (Transfer to Line 1 at City Hall Station.)
  \item \ko{다음 역에서 내리세요.} Daeum yeogeseo naeriseyo. (Get off at the next
        station.)
\end{itemize}
\end{example}

\begin{important}[Particles with \ko{타다} and \ko{내리다}]
You \ko{타다} the vehicle with \ko{을/를}: \ko{버스를 타요} (I take the bus), and you
\ko{내리다} from the vehicle or at the stop with \ko{에서}: \ko{버스에서 내려요},
\ko{명동역에서 내려요}. A very common mistake is \ko{*버스에 타요}, which some Koreans
also say, but the textbook form is \ko{를}.
\end{important}

\section{Grammar 2: \ko{에서 \til 까지} and \ko{부터 \til 까지}}

\begin{itemize}
  \item \textbf{Place}: \ko{A에서 B까지} = from A to B.
  \item \textbf{Time}: \ko{A부터 B까지} = from A until B.
\end{itemize}

\begin{example}[Distance and duration]
With \ko{걸리다} (to take time) and \ko{얼마나} (how much, how long):
\begin{itemize}
  \item \ko{집에서 학교까지 얼마나 걸려요?} Jibeseo hakgyokkaji eolmana geollyeoyo? (How
        long does it take from your home to school?)
  \item \ko{버스로 이십 분쯤 걸려요.} Beoseuro isip bunjjeum geollyeoyo. (About twenty
        minutes by bus.) --- \ko{쯤} = about.
  \item \ko{서울에서 부산까지 기차로 세 시간쯤 걸려요.} Seoureseo Busankkaji gicharo se
        sigan jjeum geollyeoyo. (From Seoul to Busan takes about three hours by train.)
  \item \ko{수업은 아홉 시부터 열두 시까지 있어요.} Sueobeun ahop sibuteo yeoldu sikkaji
        isseoyo. (Classes are from nine to twelve.)
  \item \ko{월요일부터 금요일까지 일해요.} Woryoilbuteo geumyoilkkaji ilhaeyo. (I work
        Monday to Friday.)
\end{itemize}
Note that durations use native numbers with \ko{시간} (hours) --- \ko{세 시간} --- but
Sino-Korean with \ko{분} (minutes).
\end{example}

\section{Grammar 3: Giving directions}

Directions are given with the polite command \ko{\til (으)세요} (the same form as the
honorific present) and chained with \ko{\til 아서/어서} or \ko{\til 고}. \ko{돌다} is a
\ko{ㄹ}-stem: \ko{도세요}.

\begin{center}
\begin{tikzpicture}[every node/.style={font=\scriptsize}]
  % roads
  \fill[gray!20] (-0.4,-0.4) rectangle (0.4,5.2);
  \fill[gray!20] (-4,2.6) rectangle (5,3.4);
  \fill[gray!20] (0.4,-0.4) rectangle (5,0.4);
  % crosswalk at intersection
  \foreach \y in {2.65,2.85,3.05,3.25} \draw[white, line width=1.2pt] (0.45,\y) -- (0.85,\y);
  % buildings
  \node[draw, fill=blue!10, minimum width=1.6cm, minimum height=0.8cm] at (2.6,1.6) {\ko{은행} bank};
  \node[draw, fill=green!10, minimum width=1.6cm, minimum height=0.8cm] at (2.2,4.2) {\ko{우체국} post office};
  \node[draw, fill=orange!20, minimum width=1.6cm, minimum height=0.8cm] at (-1.6,4.2) {\ko{약국} pharmacy};
  \node[draw, fill=red!15, minimum width=1.6cm, minimum height=0.8cm] at (-1.6,1.5) {\ko{지하철역}};
  \node[draw, circle, fill=yellow!60, inner sep=1pt] (you) at (0,-0.1) {\small \ko{나}};
  \node[draw, circle, fill=red!40, inner sep=2pt] at (0.7,2.25) {};
  \node[anchor=west] at (0.8,2.25) {\ko{신호등}};
  % route
  \draw[-{Stealth}, very thick, red!70!black] (0,0.25) -- (0,3) -- (2.9,3);
  \node[anchor=west, red!70!black] at (0.5,0.85) {1 \ko{쭉 가세요}};
  \node[red!70!black, anchor=west] at (3.0,3.0) {2 \ko{오른쪽으로 도세요}};
  \node[anchor=north, align=center] at (0,-0.55) {you are here};
\end{tikzpicture}
\end{center}

\begin{example}[Asking the way]
\ko{실례합니다. 우체국이 어디에 있어요?} Sillyehamnida. Uchegugi eodie isseoyo?
(Excuse me, where is the post office?)
\begin{quote}
\ko{이 길로 쭉 가세요. 사거리에서 오른쪽으로 도세요. 그리고 조금만 가면 왼쪽에 우체국이
있어요. 은행 건너편이에요.}
\end{quote}
I gil-lo jjuk gaseyo. Sageorieseo oreunjjogeuro doseyo. Geurigo jogeumman gamyeon oenjjoge
uchegugi isseoyo. Eunhaeng geonneopyeonieyo. (Go straight along this road. Turn right at
the intersection. Then if you go a little further, the post office is on your left. It's
across from the bank.) The \ko{\til (으)면} (``if'') in \ko{가면} is the next lesson's
grammar.
\end{example}

\section{Pronunciation notes}

\begin{itemize}
  \item \ko{정류장} [\ko{정뉴장}] jeongnyujang: \ko{ㄹ} after \ko{ㅇ} becomes [n].
  \item \ko{지하철로} [\ko{지하철로}] jihacheollo: \ko{ㄹㄹ} is a long [l].
  \item \ko{옆에} [\ko{여페}] yeope; \ko{앞에} [\ko{아페}] ape: liaison keeps \ko{ㅍ} as
        [p\textsuperscript{h}].
  \item \ko{몇 호선} [\ko{며토선}] myeot hoseon: aspiration across words in fast speech.
\end{itemize}

\begin{aside}
\textbf{Getting around Seoul.} The Seoul subway is numbered by line (\ko{1호선},
\ko{2호선}, \ldots) and colour; Line 2 is a loop around the city centre. A rechargeable
transit card (T-money) works on buses, subways and even taxis, and transfers between bus
and subway within a short window are free or cheap. Announcements on the subway are made
in Korean, English, and on some lines Chinese and Japanese. Seats near the doors at the
end of each car are reserved for seniors, pregnant women and people with disabilities,
and young people are expected to leave them empty even when the train is crowded ---
another everyday reflection of respect for age.
\end{aside}

\section*{Review summary}
\begin{itemize}
  \item Location: noun + \ko{앞/뒤/옆/위/아래/안/밖/사이/건너편} + \ko{에} (or
        \ko{에서} for actions).
  \item \ko{(으)로}: means (\ko{버스로}, \ko{지하철로}) and direction
        (\ko{오른쪽으로}); \ko{로} after vowels and \ko{ㄹ}.
  \item \ko{A에서 B까지} (place), \ko{A부터 B까지} (time); \ko{얼마나 걸려요?}
        --- \ko{\ldots 쯤 걸려요}.
  \item \ko{를 타다}, \ko{에서 내리다}, \ko{(으)로 갈아타다}; \ko{걸어서} = on foot.
  \item Directions: \ko{쭉 가세요}, \ko{오른쪽으로 도세요}, \ko{건너세요}.
\end{itemize}
